% ============================================================
% TITLE PAGE CONTENT — edit this file to update the title page.
% ============================================================
% This file only defines text values. It does not control layout,
% spacing, fonts, or colors — that's all in title-page.tex, which you
% shouldn't need to touch for a normal title/subtitle change.
%
% Each line below is: \newcommand{\macroName}{the text you want}
% Keep the curly braces. If your text contains a LaTeX special
% character (%, &, #, _, $) escape it with a backslash, e.g. \%.
% For a line break inside a value (like the tagline below), use \\.
%
% Every value below is a placeholder — replace with your actual
% book's title-page text.
% ============================================================

\newcommand{\giTitlePageTitle}{Your Book Title}
\newcommand{\giTitlePageCode}{COURSE-CODE}
\newcommand{\giTitlePageSubtitle}{A one-line subtitle describing this book}
\newcommand{\giTitlePageTagline}{A short tagline or description that\\expands on what this book covers}

% Which logo file to use, and how wide to render it (as a fraction of
% the page's text width — 0.42 = 42\%). logo-placeholder.png is a
% generic stand-in — replace it with a real logo file and update the
% filename here.
\newcommand{\giTitlePageLogo}{logo-placeholder.png}
\newcommand{\giTitlePageLogoWidth}{0.42}

% Optional full-bleed cover page, rendered before the title page —
% only used if cover-page.tex is added to include-before-body in
% _quarto.yml (it isn't, by default; see README's "Cover page"
% section). cover-placeholder.png is sized to A4 proportions
% (1240x1754px) so it fills the page with no distortion — a real
% cover image should be supplied at the same proportions, or larger
% at the same ratio for print quality.
\newcommand{\giCoverPageImage}{cover-placeholder.png}
